\section{Rendering on demand, observers, and quantum foundations}
\label{sec:rendering-qm}

A family of proposals holds that a resource-limited simulator would compute
(``render'') physical detail only when it is observed, and that quantum
measurement is where this shows. We first state what such a test must predict.
We then examine the most detailed proposal, by Campbell, Owhadi, Sauvageau and
Watkinson~\cite{Campbell2017on}, and summarise how existing experiments
constrain it.

\subsection{What a test must predict}
\label{sec:rendering-qm-logic}

The following argument is ours. Suppose a simulator reproduces the outcome
statistics of standard quantum mechanics (QM) exactly. Then no experiment whose
statistics QM predicts can distinguish it from a world that simply obeys QM.
This includes the delayed-choice and quantum-eraser experiments often invoked
as support. Their surprising features, such as a pattern that depends on a later
choice of measurement, are QM predictions, not anomalies. A render-on-demand
hypothesis becomes testable only when it names an observable (a fringe
visibility $V$, a conditional probability) and predicts a deviation $\Delta$ from
the QM value that exceeds the experimental uncertainty. Experiments that agree
with QM then bound $\Delta$
(Section~\ref{sec:rendering-qm-constraints}).

\subsection{The proposal of Campbell et al.\ (2017)}
\label{sec:rendering-qm-campbell}

\paragraph{Assumptions and conjecture.}
Campbell et al.\ assume that the simulating system ``is finite'' and ``seeks to
achieve low computational complexity''. By analogy with procedurally generated
games, they infer that it ``would render reality only at the moment the
corresponding information becomes available for observation by a conscious
observer''~\cite[p.~80]{Campbell2017on}. To accommodate Bell violations, they
posit that locality inside the simulation does ``not constrain the action space
of the system performing the simulation''~\cite[p.~81]{Campbell2017on}. The
simulator is therefore explicitly nonlocal with respect to simulated
spacetime. The testable conjecture is that wave or particle patterns are fixed
``by the existence and availability of the which-way data when the pattern is
observed'', not at detection~\cite[p.~82]{Campbell2017on}. ``Availability'' means
recording ``on objective media''~\cite[p.~87]{Campbell2017on}. The paper does
not say when a physical record, such as an atomic state or an emitted photon,
counts as one.

\paragraph{Proposed experiments.}
The paper proposes four schemes:
\begin{enumerate}
\item which-path detectors are left on at the slits, but the recorder is
  unplugged; success is interference. In the entangled version, the
  coincidence counter of the Kim et al.\ eraser~\cite{Kim2000delayed} is
  removed and ``D0 should display the wave pattern if the experiment is
  successful''~\cite[p.~86]{Campbell2017on};
\item which-path and screen data go to separate USB drives, and the which-path
  drive is irreversibly destroyed with probability $1/2$; success is
  interference only where the partner drive was
  destroyed~\cite[pp.~87--88]{Campbell2017on};
\item in a delayed-choice eraser, the signal position $X$ predicts the later
  random which-path/erasure outcome $R$ via
  $P[R{=}1|X{\approx}x]\approx[1+2\cos^2(\pi x/a)]^{-1}$ with
  $a=\lambda L/d$~\cite[Eq.~(2)]{Campbell2017on};
\item a thought experiment with a manual switch and delays of about 60\,s. A
  total-variation argument excludes $X$ being drawn from the wave distribution
  with the switch off and from the particle distribution with it
  on~\cite[Eqs.~(3)--(4)]{Campbell2017on}.
\end{enumerate}
The authors ``cannot predict the outcome'' of scheme~4 but prove it ``will be
new''~\cite[p.~83]{Campbell2017on}. No scheme carries a numerical prediction.
The only tolerance given is a noise allowance $\delta(I)<0.9$ in
scheme~4~\cite[p.~94]{Campbell2017on}.

\paragraph{Do the schemes predict deviations from QM?}
They do, although the paper does not compare its success criteria with QM. The
comparison below is ours and draws on the cited literature. In QM, which-path
information carried by any system suppresses interference, to the degree that the
carrier's states are distinguishable, whether or not anyone reads it. The information ``is still present in the universe, independent of
whether an observer takes note of it or not'', and ignoring its carrier leaves
a mixed state ``which cannot show an interference
pattern''~\cite[Sec.~II.E]{Ma2016delayed}. Schemes~1 and~2 therefore demand
visibilities of order unity where QM bounds $V$ by the which-path parameter $I$,
through $I^2+V^2\le1$~\cite{Englert1996fringe,Ma2013quantum}. In QM,
destroying a classical copy of the record does not restore interference. Interference
reappears only conditionally, through a suitable measurement on the quantum
carrier of the path information~\cite{Ma2016delayed,Kastner2019delayed}.

In the Kim et al.\ eraser, D0 coincidences with the two erasing detectors give
fringes and antifringes ($R_{01}\propto\cos^2$, $R_{02}\propto\sin^2$). The
which-path coincidences give none~\cite[Eq.~(10)]{Kim2000delayed}. The unsorted
D0 data are thus fringe-free, ``just noise'' in Kastner's
words~\cite{Kastner2019delayed}. Removing the coincidence counter (scheme~1)
therefore cannot produce fringes in QM. The erasure subensemble (pooling the two erasing detectors) and the
which-path subensemble are both fringe-free, so in QM $X$ and $R$ are independent and
$P(R{=}1|X)=1/2$. Eq.~(2)
instead assumes a single-detector fringe for $R=0$. It departs from QM by up to
$1/2$ in probability at dark fringes, where it gives $\approx1$, and by $1/6$ at
bright fringes, where it gives $1/3$. Scheme~4 assumes ``an interference pattern
at D0'' in its configurations (a)--(c)~\cite[p.~91]{Campbell2017on}. That is
also not a QM prediction for an entangled-pair eraser, where fringes appear
only after conditioning on the partner
photon~\cite{Ma2013quantum,Kastner2019delayed}. The paper classifies the
possible outcomes as a ``discontinuity in the rendering reality'', rendering at
observation, or a paradox~\cite[p.~91]{Campbell2017on}. Under this scheme, the
QM result (no fringes in unconditioned D0 data in any configuration) counts as
a ``discontinuity''. In configuration (d) it coincides with the outcome the
authors call ``a strong indicator that this reality is
simulated''~\cite[p.~92]{Campbell2017on}. As written, the protocol has no
outcome that counts against the hypothesis.

\paragraph{Status.}
The non-profit CUSAC, founded by Campbell, lists five experiments ``that CUSAC
intends to perform''. It names leads in California (F.\ Khoshnoud) and Canada
(D.\ Chartrand) but posts no data~\cite{CUSAC2024experiments}. A paper
co-authored by Khoshnoud, which acknowledges Campbell and Chartrand, reports a
teaching-lab single-photon double slit with $|V|=0.65\pm0.22$. It also reports a
polarisation eraser in which orthogonal polarisers on the slits removed the
fringes without any measurement: ``we did not need to measure the path
information''~\cite{Luo2024young}. That paper is not presented as a test of the
simulation hypothesis, and a polarisation marker is not the unplugged recorder
of scheme~1. As of September 2026 we found no citable report of any of the four
schemes having been performed.

\subsection{What existing experiments constrain}
\label{sec:rendering-qm-constraints}

Wheeler proposed choosing the measurement after the photon has entered the
interferometer~\cite{Wheeler1978past}, and later described a cosmological
version~\cite{Wheeler1983law}. His criterion for a completed phenomenon was an
``irreversible act of amplification'', not conscious
awareness~\cite[Sec.~II.D]{Ma2016delayed}. The quantum
eraser~\cite{Scully1982quantum} extended the scheme to entangled pairs.
Table~\ref{tab:rendering-qm-exps} lists realisations that test the two
questions a rendering hypothesis raises: does the timing of the choice matter,
and does an unread record matter? All agree with QM within their quoted
uncertainties. The review by Ma, Kofler and Zeilinger concludes that ``no
physical interactions or signals, let alone into the past, are necessary to
explain the experimental results''~\cite[Sec.~VI]{Ma2016delayed}.

\begin{table}[t]
\centering
\small
\caption{Experiments that bound render-on-demand deviations ($V$: fringe
visibility; $I$: which-path parameter; uncertainties as quoted, $\pm1\sigma$
where stated).}
\label{tab:rendering-qm-exps}
\begin{tabular}{p{0.25\linewidth}p{0.34\linewidth}p{0.33\linewidth}}
\hline
Experiment & Varied, or left unread & Result \\
\hline
Jacques et al.~\cite{Jacques2007experimental} &
Open/closed choice by QRNG, space-like separated from entry (48\,m) &
Closed: $V=94\%$; open: no fringes, port probabilities $0.50\pm0.01$, $I>0.99$ \\
Kim et al.~\cite{Kim2000delayed} &
Idler registered $\geq8$\,ns after the signal &
Fringes/antifringes with D1/D2; none with D3/D4 \\
Ma et al.\ 2013~\cite{Ma2013quantum} &
Erasure choice space-like separated (55\,m, 144\,km), or $\approx450\,\mu$s later &
$I=0.955(7)$, no fringes; erasure: $V=0.951(18)$, $I=0.077(22)$ \\
Ma et al.\ 2012~\cite{Ma2012experimental} &
Swapping choice 14--313\,ns after the other photons were registered &
Fidelity $0.681\pm0.034$ vs $0.421\pm0.029$ \\
D\"urr et al.~\cite{Durr1998origin} (via~\cite{Ma2016delayed}) &
Path stored in atomic internal states, not read out &
Fringes vanish; momentum kick $10^{-4}$ of fringe period \\
Hackerm\"uller et al.~\cite{Hackermuller2004decoherence} &
Path carried off by undetected thermal photons &
$V=47,29,7,0\%$ at 0, 3, 6, 10.5\,W; matches decoherence theory \\
Hornberger et al.~\cite{Hornberger2003collisional} &
Path left in colliding gas molecules &
Pressure-dependent loss of $V$, as predicted \\
\hline
\end{tabular}
\end{table}

Two limits apply. First, the delayed-choice experiments vary the timing of the
choice. None of them destroyed a classical record, or varied whether a person
viewed it, as schemes~1--2 require. Second, the decoherence and internal-state
experiments show interference lost to unread records. D\"urr et al.\ put this
as ``the mere fact that which-path information is stored in the detector and
could be read out already destroys the interference
pattern''~\cite[Sec.~IV.B]{Ma2016delayed}. Because ``availability'' is not
operationally defined, however, a proponent could call such records
``available''. On our reading, any deviation tied to \emph{when} a choice is
made is bounded at the few-per-cent level of the best photonic $V$ and $I$
uncertainties. A test of the unread-record variant requires fixing the
availability criterion before the data are taken.

\subsection{Other observer-dependent and ``glitch'' proposals}
\label{sec:rendering-qm-others}

Whitworth argued, in a technical report, that a virtual world ``would likewise
be expected to be calculated only on demand''~\cite{Whitworth2008physical}.
Arvan proposed a ``peer-to-peer simulation hypothesis'' for quantum
phenomena~\cite{Arvan2014unified}. Irwin, Amaral and Chester's
``self-simulation'' interpretation~\cite{Irwin2020self} asserts that it
``became more widely agreed that it is consciousness'', rather than interaction
with a detector, that changes the double-slit pattern. This is at odds with the
review literature above~\cite{Ma2016delayed}. As evidence they cite Radin et
al.'s attention experiments~\cite{Radin2016psychophysical}. An independent
reanalysis of the 2016 dataset found shifts in the predicted direction but ``not
deemed significant ($p>0.05$)'', and concluded that it ``does not contain
evidence of mind-matter interaction''~\cite{Tremblay2019independent}.
Yampolskiy's ``How to hack the simulation?'', published as ``How to Escape From
the Simulation''~\cite{Yampolskiy2023how}, does not evaluate evidence. It
proposes that ``every novel QM experiment can be seen as an attempt at hacking
the simulation'', and it surveys proposals to overload the simulator with
computationally intense activity. A reviewer comment printed with the paper
argues that if simulators cannot optimise through layers of abstraction, a
full-universe simulation is more likely. That would make it futile to expand
``the human-observed-in-detail cone''~\cite{Yampolskiy2023how}. None of
these works gives a quantitative prediction that differs from QM. By contrast, Chalmers and McQueen give explicit dynamics for
consciousness-triggered collapse and report that simple versions are
``falsified by the quantum Zeno effect''. More complex versions remain
compatible with current data and are in principle testable with quantum
computers~\cite{Chalmers2022consciousness}. Their model is not a simulation
proposal, but it shows the kind of formulation the render-on-demand literature
lacks.

\subsection{Nonlocality, superdeterminism and freedom of choice}
\label{sec:rendering-qm-bell}

Loophole-free Bell tests constrain any simulator that works from local data.
Hensen et al.\ found $S=2.42\pm0.20$ in 245 trials at 1.3\,km ($p=0.039$)
\cite{Hensen2015loophole}. Giustina et al.\ reported $p\le3.74\times10^{-31}$
\cite{Giustina2015significant}. Shalm et al.\ reported $p$ down to
$5.9\times10^{-9}$, or $2.3\times10^{-7}$ after adjusting for setting
predictability~\cite{Shalm2015strong}. A simulator matching these results must
therefore act nonlocally with respect to simulated spacetime, as Campbell et
al.\ posit, or correlate its hidden state with the settings. The latter is a
violation of Statistical Independence, i.e.\
superdeterminism~\cite{Hossenfelder2020rethinking}. Campbell et al.'s ``VR
engine reacting to the intention of the experimentalist''~\cite[p.~94]{Campbell2017on}
can be read as falling in this class.

In-universe versions of the setting-correlation option have been narrowed.
The BIG Bell Test used 97{,}347{,}490 choices from about 100{,}000
participants~\cite{BIGBellTest2018challenging}. Cosmic Bell tests used Milky
Way stars, with violations of $7.31\sigma$ and $11.93\sigma$ quoted as lower
bounds, pushing the latest local-realist intervention back by
$\sim600$\,yr~\cite{Handsteiner2017cosmic}. With quasars they reached $9.3\sigma$
and pushed it back to at least $\sim7.8$\,Gyr ago, covering 96\% of the
past-light-cone space-time volume~\cite{Rauch2018cosmic}. The cosmic tests assume
fair sampling, and all of these tests constrain \emph{local} mechanisms. Hossenfelder and Palmer
argue that they ``do not -- cannot -- rule out Superdeterminism'', and that
``one cannot prove freedom of choice by assuming freedom of
choice''~\cite[Sec.~4.3]{Hossenfelder2020rethinking}. Neither kind of simulator
above is constrained by these bounds. A nonlocal one is not, because Bell
tests exclude only local models. One that sets settings and outcomes jointly
from outside spacetime is not either, because it is empirically a
superdeterministic model. It therefore inherits superdeterminism's testability
problem: the signature Hossenfelder and Palmer propose, Born-rule-violating
time-correlations for nearly identical preparations, needs a concrete model to
fix how similar the preparations must be~\cite[Sec.~6]{Hossenfelder2020rethinking}.
't~Hooft uses a computer image himself: a God who must ``run a universe''
prefers deterministic rules on a classical computer. He sees ``no objections
against super-determinism''~\cite{tHooft2019free}. His cellular automaton
interpretation predicts that quantum computers cannot outperform a classical
computer with one memory site per Planck volume updating once per Planck time.
He states that exceeding this bound would falsify the
theory~\cite[Sec.~5.8]{tHooft2016cellular}. That is a quantified,
resource-based prediction of the kind the render-on-demand proposals reviewed
here do not yet provide.
